\chapter{Generación de señales, reconstrucción y procesamiento}
\section{GNU Radio}
Para generar la señal ISDB-Tb se utilizaron los bloques de procesamiento publicados en~\cite{grisdbt_repo}. A partir de estos bloques se construyeron los flujos de transmisión, recepción y remodulación mostrados en las Figuras~\ref{fig:isdbt_iluminador}--\ref{fig:tx_remod}.

\figura{\Tfig}{./gnu_radio/tx.png}{Generación de la señal ISDB-Tb del iluminador de oportunidad.}{}{\label{fig:isdbt_iluminador}}

\figura{\Tfig}{./gnu_radio/rx_A.png}{Recepción y demodulación de la referencia ruidosa (Capa A).}{}{\label{fig:rx_A}}
\figura{\Tfig}{./gnu_radio/rx_B.png}{Recepción y demodulación de la referencia ruidosa (Capa B).}{}{\label{fig:rx_B}}

\figura{\Tfig}{./gnu_radio/tx_remod.png}{Remodulación para obtener referencia reconstruida.}{}{\label{fig:tx_remod}}

La Figura~\ref{fig:isdbt_iluminador} muestra el flujo utilizado para generar la señal ISDB-Tb asociada al iluminador de oportunidad. Luego, las Figuras~\ref{fig:rx_A} y~\ref{fig:rx_B} presentan los flujos de recepción correspondientes a las capas A y B, respectivamente. Finalmente, la Figura~\ref{fig:tx_remod} muestra el flujo de remodulación empleado para obtener la referencia reconstruida, la cual será comparada con la referencia ruidosa en la cadena de procesamiento.

\newpage
\section{Implementación de las cadenas de procesamiento}
\subsection{Generación de canal de vigilancia y referencia}
El Listado~\ref{lst:gen_signals} presenta la función encargada de generar las señales utilizadas en la simulación. Esta función agrega ruido complejo a la señal ISDB-Tb de referencia y configura una escena de radar pasivo para obtener la referencia ruidosa y el canal de vigilancia que luego se usaran en la cadena de procesamiento.\\\vspace{5pt}
\begin{lstlisting}[
    style=pythoncompact,
    caption={Función para generar las señales de referencia y vigilancia simuladas.},
    label={lst:gen_signals}
]
def gen_signals(snr_db, reference, echo_power_db, samples=N_SAMPLES):
    N = len(reference)
    E_ref = np.mean(np.abs(reference) ** 2)
    noise_power = E_ref / utils.from_db(snr_db)
    noise = (
        np.sqrt(noise_power)
        * (np.random.randn(N) + 1j * np.random.randn(N))
    ).astype(np.complex64)

    noisy_ref = (reference + noise).astype(np.complex64)
    noisy_ref.tofile(DATA_DIR / "gr_files" / "isdbt_noisy.cfile")

    pr_config = PassiveRadarChainConfig(
        input=InputConfig(
            N=samples, fs=8126984.0, f_c=700e6, use_simulated_data=True
        ),
        simulation=SimulationConfig(
            transmitter_position=[0.0, 0.0],
            radar_position=[50.0, 500.0],
            direct_signal=True,
            clutter=ClutterConfig(
                N_CLUTT=10,
                clutter_rcs_min_db=-3.0,
                clutter_rcs_max_db=-5.0,
                rand_clutter=True,
                clutter_limits=[-100, 2000, 5, 150],
            ),
            echo=EchoConfig(
                V_b=[10.0, 100.0],
                rand_target=False,
                target_rcs_db=echo_power_db,
                target_position=[2500.0, 120.0],
            ),
        ),
        channel=ChannelConfig(
            enable=False,
        ))
    
    noise_surv = (
        np.sqrt(noise_power)
        * (np.random.randn(samples) + 1j * np.random.randn(samples))
    )

    pr_chain = PassiveRadarChain(pr_config)
    pr_chain.simulate_inputs(reference[:samples])    
    
    input_state= pr_chain.get_state(stage= "inputs", copy_state=False)
    sim_state = pr_chain.get_state(stage="simulation", copy_state=False)

    surv = input_state.surveillance + noise_surv
    fd, target_p = sim_state.doppler_hz, sim_state.target_position

    return noisy_ref[0:samples], surv, fd, target_p
\end{lstlisting}
\newpage
\subsection{Ejecución de la cadena de procesamiento}
El Listado~\ref{lst:run_config} muestra la función encargada de ejecutar la cadena de procesamiento. Esta función selecciona la referencia en función de si se utiliza la reconstrucción o no, carga el canal de vigilancia y ejecuta progresivamente las etapas de cálculo de la CAF, filtrado de clutter, aplicación de ventanas y detección CFAR. Además, permite guardar las figuras, la configuración y el estado final de la cadena.\\\vspace{5pt}

\begin{lstlisting}[
    style=pythoncompact,
    caption={Funcion encargada de ejecutar la cadena de procesamiento (con y sin reconstrucción).},
    label={lst:run_config}
]
def run_config(remod, pass_parameters, beta, cfar, samples=N_SAMPELS, save=False) -> None:
    REMOD_TITLE = "_remod" if remod else "_no_remod"
    RECONSTRUCTOR_TITLE = "con reconstrucción" if remod else "sin reconstrucción"
    if pass_parameters:
        chain = PassiveRadarChain(
            config=build_config(cfar=cfar, beta=beta,samples=samples, save=save), verbose=True
        )
    else:
        chain = PassiveRadarChain(config=build_config(), verbose=True)

    signals = np.load(DATA_DIR / "pr_signals" / "signals.npz")
    
    surv = signals['surv']
    if remod:
        ref = signals['remod']
    else:
        ref = signals['noisy']

    fig1, ax1 = plt.subplots()
    utils.plot_psd(x= ref, fs= 8126984.0, ax=ax1, n_samples=N_SAMPELS, freq_in_khz=True)
    ax1.yaxis.set_major_locator(MaxNLocator(nbins=15))
    ax1.set_title(f"Señal de referencia {RECONSTRUCTOR_TITLE}")
    if save:
        fig1.savefig(FIG_DIR / f"ref{REMOD_TITLE}.png", dpi=300, bbox_inches="tight")
        
    chain.set_inputs(reference= ref, surveillance= surv)
    # CAF Sin Filtro CF ni ventanas
    chain.run_until(stage="caf")
    chain.plot_caf(
        filename=f"caf{REMOD_TITLE}.png",
        title=f"CAF {RECONSTRUCTOR_TITLE}",
    )

    # CAF con Filtro CF y sin ventanas
    chain.update_filter_config(enabled=True)
    chain.run(start_from="filter", stop_at="caf")
    chain.plot_caf(
        filename=f"filtered_caf{REMOD_TITLE}.png",
        title=f"CAF {RECONSTRUCTOR_TITLE} (CF)",
    )

    # CAF con Filtro CF y ventanas
    chain.update_window_config(enabled=True)
    chain.run(start_from="window", stop_at="detect")
    chain.plot_caf(
        filename=f"filtered_w_caf{REMOD_TITLE}.png",
        title=f"CAF {RECONSTRUCTOR_TITLE} (CF + Ventanas)",
    )

    chain.plot_detections(
        filename=f"filtered_w_detections{REMOD_TITLE}.png",
        title=f"Detecciones CAF {RECONSTRUCTOR_TITLE} (CF + Ventanas) ",
    )
    chain.save_config(filename=f"config{REMOD_TITLE}")
    chain.save_state(filename=f"state{REMOD_TITLE}")

\end{lstlisting}
\newpage
\subsubsection{Testeo de cadenas}
El Listado~\ref{lst:test_full_chain} contiene la función principal utilizada para coordinar la prueba completa. Esta función permite ejecutar los archivos \texttt{.py} generados a partir de los flujos de GNU Radio, cargar las señales generadas, construir la referencia ruidosa y almacenar todos los datos necesarios en un archivo \texttt{.npz}. Finalmente, llama a la función de procesamiento para ejecutar ambos casos de análisis: con referencia reconstruida y con referencia ruidosa.\\\vspace{5pt}
\begin{lstlisting}[
    style=pythoncompact,
    caption={Función para generar señales, almacenar datos y ejecutar los casos con y sin reconstrucción.},
    label={lst:test_full_chain}
]
def test_full_chain(gen_isdtb_signals:bool = False):
    if gen_isdtb_signals:
        subprocess.run([sys.executable, DATA_DIR/"gr_files"/"tx_tesis.py"])
    reference = np.fromfile(DATA_DIR / "gr_files" / "isdtb_clean.cfile", dtype=np.complex64)

    snr = 15
    noisy_ref, surv, fd, target_p = gen_signals(snr_db = snr, reference= reference, echo_power_db= -30.0, samples=N_SAMPLES)
    

    
    if gen_isdtb_signals:
        subprocess.run([sys.executable, DATA_DIR/"gr_files"/"rx_tesis_A.py"])
        subprocess.run([sys.executable, DATA_DIR/"gr_files"/"rx_tesis_B.py"])
        subprocess.run([sys.executable, DATA_DIR/"gr_files"/"tx_tesis_remod.py"])
    
    reference_remod = np.fromfile(DATA_DIR / "gr_files" / "isdbt_remod.cfile", dtype=np.complex64)

    np.savez(
    DATA_DIR / "pr_signals"/"signals.npz",
    remod=reference_remod[:N_SAMPLES],
    noisy=noisy_ref,
    surv= surv,
    target_doppler = fd,
    target_p = target_p
    )
    
    run_config(remod= True, pass_parameters= True, beta=(40.0, 20.0), cfar=(1, (128, 512)),samples=N_SAMPLES, save= True)
    run_config(remod= False, pass_parameters= True, beta=(40.0, 20.0), cfar=(1, (128,512)),samples=N_SAMPLES, save= True)
\end{lstlisting}